\section{Introduction}
Attribution graphs built from sparse replacement models are becoming standard tools for explaining how language models compute~\cite{ameisen2025circuit,lindsey2025biology}. Cross-layer transcoders (CLTs) let a feature that reads the residual stream at layer $\ell$ write to every later MLP output. This handles features that are computed across several layers and gives graphs with many fewer nodes~\cite{ameisen2025circuit}. The same freedom has a cost. A CLT can explain a deep, multi-step computation with features that jump straight from the inputs to the output. The resulting graph matches the model's behaviour while hiding its mechanism. On a two-layer residual MLP computing $(a\oplus b)\wedge(c\oplus d)$, recent work showed that CLTs do exactly this~\cite{clt_unfaithful2026}.

That study used one depth and an attention-free model. Three questions remain open. (i) Does the collapse \emph{grow} with the true number of sequential steps? (ii) Does it hold when CLTs are compared with per-layer transcoders (PLTs) at a strictly identical budget, rather than at the budgets each would naturally be trained with? (iii) When a graph is shallow, is the model itself shallow, or is the graph wrong? Answering (iii) needs a ground truth whose depth is \emph{certified}, not assumed.

\paragraph{Contributions.}
\begin{enumerate}
\item \textbf{A certified depth sweep.} We train attention transformers with strict interchange intervention training (SIIT)~\cite{gupta2024interpbench}, adding a bypass constraint and reserved site coordinates, on Boolean tasks of depth $D=1..5$. Every model must pass interchange, strictness, bypass and circuit-only necessity tests. We show that a standard 2-hop design with a final AND has nonlinear depth 1, because the unembedding performs the last hop.
\item \textbf{A depth-collapse metric.} Our metric is the effective graph depth $\Dfeat$: the average number of feature hops on the paths from the root inputs to the output, weighted by their share of the contribution. It comes in a cancellation-aware form and in the Circuit-Tracing absolute-value form. The depth-collapse ratio is $\Dfeat/D$.
\item \textbf{A matched comparison.} PLTs and CLTs get the same number of features per layer, the same exact $L_0$ (TopK), the same training, and the same base models (a paired design). We add a parameter-matched PLT control and a sparsity sweep.
\item \textbf{Claims versus causes.} We separate \emph{decodable} from \emph{causally used} intermediates, and \emph{reconstruction} from \emph{validated} mechanisms, using route-freezing interventions in the replacement model. We validate the linearised graph against real ablations.
\item \textbf{Beyond toys.} We evaluate two-hop factual recall in GPT-2, whose intermediate entity is known~\cite{lindsey2025biology}, with one-hop controls, feature ablations and graph--intervention comparisons. InterpBench remains a proposed extension and was not run.
\end{enumerate}
